\documentclass[11pt,a4paper]{book}


\input{../1ere-preambule}

\begin{document}
\graphicspath{{images/}}
%\setcounter{chapter}{1}
\chapter{Automatismes}
\thispagestyle{fancy}

\section{proportion et pourcentage}

\subsection{Calculer, appliquer, exprimer une proportion sous différentes formes}

\begin{bmethode}
    \grasInMet{Enoncé :}
    
    On réalise un sondage auprès de 400 personnes concernant les mesures prises par le gouvernement.
    \begin{enumMet}
        \item 94 de ces 400 personnes ont affirmés être pleinement satisfaite des mesures prises par le gouvernement.
        
        Déterminer la proportion de personnes pleinement satisfaite des mesures prises par le gouvernement, puis l'exprimer sous forme de pourcentage.
        \item La proportion des personnes non satisfaites est $\dfrac{3}{8}$.
        
        Combien de personnes ont affirmés ne pas être satisfaites des mesures prises par le gouvernement?
    \end{enumMet}
    \grasInMet{Réponse :}
    
    \begin{enumMet}
        \item Le nombre de personnes interrogés est $n_E=400$. Parmi ceux-ci, le nombre de ceux satisfait est $n_S=94$.
        
        La proportion de personnes pleinement satisfaite des mesures prises par le gouvernement est $p=\dfrac{n_S}{n_E}=\dfrac{94}{400}=0,235=23,5\%$.
        \item Soit $n_N$ le nombre de personnes non satisfaites.
        
        On a $\dfrac{n_N}{n_E}=\dfrac{3}{8}$, soit $\dfrac{n_N}{400}=\dfrac{3}{8}$.
        On a donc $n_N=\dfrac{400\times 3}{8}=150$.
        
        Il y a donc 150 personnes non satisfaites des mesures prises par le gouvernement.
    \end{enumMet}
\end{bmethode}

\subsection{Calculer la proportion d'une proportion}

\begin{bmethode}
    \grasInMet{Enoncé :}
    
    Un maraîcher vend des légumes en direct à la ferme et sur des marchés mais aussi dans des supermarchés locaux. Au cours du mois de Juin, il a vendu 78\% de sa production en direct, et parmi ces légumes, 65\% ont été vendu à la ferme.
    
    Quelle proportion de sa production a été vendu directement à la ferme?
    
    \hfill 
    
    
    \grasInMet{Réponse :}
    
    La proportion $p_1$ de légumes vendus en direct est $p_1=0,78$
    
    La proportion $p_2$ de légumes vendus à la ferme parmi ceux vendus en direct est $p_2=0,65$
    
    On calcule : $p=p_1\times p_2=0,78\times0,65=0,507=50,7\%$
    
    La proportion de sa production vendu directement à la ferme a été de $50,7\%$
  
\end{bmethode}

\subsection{Passer du pourcentage d'évolution au coefficient multiplicateur}

\begin{bmethode}
    \grasInMet{Enoncé :}
    
    Le nombre de naissances dans un département français entre 2010 et 2019 a diminué de $12\%$
    \begin{enumMet}
        \item Par quel coefficient a été multiplié le nombre de naissances entre 2010 et 2019.
        \item Le nombre d'habitants de ce département a, quant à lui, été multiplié par $1,03$.
        
        A quel taux d'évolution du nombre d'habitants de ce département ce coefficient multiplicateur correspond-il?
    \end{enumMet}
    \grasInMet{Réponse :}
    
    \begin{enumMet}
        \item Le taux d'évolution du nombre de naissances est $t=-0,12$.
        
        Diminuer une quantité de $12\%$ signifie la multiplier par $c=1-\dfrac{12}{100}=0,88$.
        
        Par conséquent, le nombre de naissances dans ce département a été multiplié par $0,88$.
        \item Le nombre d'habitants de ce département a, quant à lui, été multiplié par $c=1,03$.
        
        Le taux d'évolution associé à ce coefficient multiplicateur est:$t=c-1=0,03$.
        
        Ainsi le nombre d'habitants du département a augmenté de $3\%$
    \end{enumMet}
\end{bmethode}

\section{Evolutions et variations}

\subsection{Appliquer un taux d'évolution pour calculer une valeur finale}

\begin{bmethode}
    \grasInMet{Enoncé :}
    
    En 2016, 88200 personnes ont visité un parc d'attractions
    \begin{enumMet}
        \item Ce nombre de visiteurs a augmenté de $12,5\%$ en 2017.
        
        Quel a été le nombre de visiteurs en 2017.
        \item Il y a eu $112500$ visiteurs en 2018, puis une baisse de la fréquentation de $8\%$ en 2019. Quel a été le nombre de visiteurs en 2019? 
    \end{enumMet}
    \grasInMet{Réponse :}
    
    \begin{enumMet}
        \item Augmenter une quantité de $12,5\%$ signifie la multipler par $c=1+\dfrac{12,5}{100}=1,125$.
        
        Or $88200\times1,125=99225$. Le nombre de visiteurs a été de $99225$ personnes en 2017.
        \item Diminuer une quantité de $8\%$ signifie la multiplier par $c=1-\dfrac{8}{100}=0,92$.
        
        Or $112500\times0,92=103500$. Le nombre de visiteurs a été de $103500$ personnes en 2019.
    \end{enumMet}
\end{bmethode}

\subsection{Appliquer un taux d'évolution pour calculer une valeur initiale}

\begin{bmethode}
    \grasInMet{Enoncé :}
    
    Le prix du ticket de métro d'une grande ville augmente de $4\%$ au $1^{\text{er}}$ janvier 2020.
    
    Le ticket de métro au tarif réduit coûte $1,56$ \euro{} après son augmentation.
    
    Combien coûtait-il avant l'augmentation?
    
    \begin{minipage}[c]{0.65\linewidth}
        \grasInMet{Réponse :}
        
        Augmenter une quantité de $4\%$ revient à la multiplier par $1,04$.
        
        Diviser par $1,04$ permet de compenser une multiplication par $1,04$.
        
        Le ticket de métro au tarif réduit coûte $1,56$ \euro{}. Donc son prix avant augmentation est donné par le quotient $\dfrac{1,56}{1,04}$.
        
        Or $\dfrac{1,56}{1,04}=1,5$. Ainsi le ticket de métro valait $1,50$ \euro{} avant le $1^{\text{er}}$ janvier 2020.
    \end{minipage} \hfill
    \begin{minipage}[c]{0.30\linewidth}
        \psset{xunit=1cm,yunit=1cm,algebraic=true,dimen=middle,dotstyle=o,dotsize=5pt 0,linewidth=1.6pt,arrowsize=3pt 2,arrowinset=0.25}
        \begin{pspicture*}(5,7)(9.5,12)
        \rput[B](7.1,10.35){\red $\times 1,04$}
        \rput[B](7.1,8.05){\color{ForestGreen} $\div 1,04$}
        \rput[B](6,8.9){\large ?}
        \rput[B](8.2,8.9){\large $1,56$}
        \psline[linewidth=1.5pt,linearc=1.5,arrows=->,arrowscale=1.3](6,9.7)(7.1,10.7)(8.2,9.7)
        \psline[linewidth=1.5pt,linearc=1.5,arrows=<-,arrowscale=1.3](6,8.3)(7.1,7.3)(8.2,8.3)
        \pscircle[linewidth=1.5pt](6.,9.){0.6}
        \pscircle[linewidth=1.5pt](8.2,9.){0.6}
        \end{pspicture*}
    \end{minipage}
     
\end{bmethode}

\subsection{Calculer un taux d'évolution et l'exprimer en pourcentage}

\begin{bmethode}
    \grasInMet{Enoncé :}
    
    Le prix d'un kilogramme d'abricots passe de $2,50$ \euro{} à $1,80$ \euro{} en début de saison.
    \begin{enumMet}
        \item Déterminer le pourcentage d'évolution du prix du kilogramme d'abricots en début de saison.
        \item En fin de saison, le prix d'un kilogramme passe de $1,80$ \euro{} à $2,61$ \euro{}.
        
        Déterminer le pourcentage d'augmentation du prix d'un kilogramme d'abricots en fin de saison
    \end{enumMet}
    \grasInMet{Réponse :}
    
    \begin{enumMet}
        \item On a $Q_1=2,50$ et $Q_2=1,80$. Le taux d'évolution $t$ de la valeur $Q_1$ à la valeur $Q_2$ est:
        
        $t=\dfrac{Q_2-Q_1}{Q_1}=\dfrac{1,80-2,50}{2,50}=\dfrac{-0,70}{2,50}=-0,28=-28\%$.
        
        Le prix du kilogramme d'abricots a diminué de $28\%$ en début de saison.
        \item On a $Q_1=1,81$ et $Q_2=2,61$. Le taux d'évolution $t$ correspondant est:
        
        $t=\dfrac{Q_2-Q_1}{Q_1}=\dfrac{2,61-1,80}{1,80}=\dfrac{0,60}{1,80}=0,45=45\%$.
        
        Le prix du kilogramme d'abricots a augmenté de $45\%$ en fin de saison.
    \end{enumMet}
\end{bmethode}

\subsection{Calculer le taux d'évolution de plusieurs évolutions successives}

\begin{bmethode}
    \grasInMet{Enoncé :}
    
    Le cours d'une action au $1^{\text{er}}$ février 2020 a augmenter de $12\%$ par rapport à son cours au $1^{\text{er}}$ janvier 2020. Il augmente ensuite de $8\%$ entre le $1^{\text{er}}$ février 2020 et le $1^{\text{er}}$ mars 2020.
    
    \begin{enumMet}
        \item Déterminer le taux d'évolution global du cours de cette action entre le $1^{\text{er}}$ janvier et le $1^{\text{er}}$ mars 2020.
        \item Le cours de cette action diminue ensuite entre le $1^{\text{er}}$ mars et le $1^{\text{er}}$ avril 2020 de $25\%$, puis augmente de $14\%$ entre le $1^{\text{er}}$ avril et le $1^{\text{er}}$ mai. Déterminer le pourcentage d'évolution du cours de l'action entre le $1^{\text{er}}$ février et le $1^{\text{er}}$ mai 2020.
    \end{enumMet}
    \grasInMet{Réponse :}
    
    \begin{enumMet}
        \item Augmenter une quantité de $12\%$ revient à la multiplier par $1+\dfrac{12}{100}$, soit $1,12$.
        
        Le coefficient multiplicateur global pour ces deux évolutions successives est donc $c=1,12\times1,08=1,2096$. Le taux d'évolution global du cours de l'action est donc $1,2096-1$, soit $0,2096$. Ainsi le pourcentage d'augmentation de ce cours est $20,96\%$. 
        \item Entre le $1^{\text{er}}$ février et le $1^{\text{er}}$ mai, il y a trois évolutions successives de coefficients multiplicateurs $1,08$, $1-0,25=0,75$ et $1+0,14=1,14$. Le coefficient multiplicateur global pour ces trois évolutions successives est donc $c'=1,08\times0,75\times1,14=0,9234$.
        
        $c'-1=-0,0766$, donc le cours de cette action a diminué de $7,66\%$ entre le $1^{\text{er}}$ février et le $1^{\text{er}}$ mai 2020. 
    \end{enumMet}
\end{bmethode}

\vspace*{-0.8cm}
\subsection{Calculer un taux d'évolution réciproque}

\begin{bmethode}
    \grasInMet{Enoncé :}
    
    La population de bouquetins dans un parc national a diminué en 2019 de $18\%$.
    
    Déterminer, à $0,1\%$ près, le pourcentage d'augmentation nécessaire de cette population en 2020 afin que la population de bouquetins retrouve son niveau initial.
    
        \begin{minipage}[c]{0.65\linewidth}
        \grasInMet{Réponse :}
        
        Diminuer une quantité de $18\%$ revient à la multiplier par $c=1-0,18=0,82$.
        Le coefficient multiplicateur de l'évolution réciproque est $c'$ tel que $c\times c'=1$, soit $c'=\dfrac{1}{c}=\dfrac{1}{0,82}$ et $c'\approx 1,22$
        
        Soit $t'$ le taux d'augmentation cherché. Alors $t'=c'-1$; donc $t'\approx 0,22$. 
        
        On en déduit que le pourcentage d'augmentation de cette population doit être de $22\%$, à $0,1\%$ près, pour revenir au niveau initial de 2018.
    \end{minipage} \hfill
    \begin{minipage}[c]{0.30\linewidth}
        \psset{xunit=1cm,yunit=1cm,algebraic=true,dimen=middle,dotstyle=o,dotsize=5pt 0,linewidth=1.6pt,arrowsize=3pt 2,arrowinset=0.25}
        \begin{pspicture*}(5,7)(9.5,12)
            \rput[B](7.1,10.35){\red $\times 0,82$}
            \rput[B](7.1,10.8){\red Diminuer de $18\%$}
            \rput[B](7.1,8.05){\color{ForestGreen} $\div 0,82$}
            \rput[B](6,8.9){\large $B_1$}
            \rput[B](8.2,8.9){\large $B_2$}
            \psline[linewidth=1.5pt,linearc=1.5,arrows=->,arrowscale=1.3](6,9.7)(7.1,10.7)(8.2,9.7)
            \psline[linewidth=1.5pt,linearc=1.5,arrows=<-,arrowscale=1.3](6,8.3)(7.1,7.3)(8.2,8.3)
            \pscircle[linewidth=1.5pt](6.,9.){0.6}
            \pscircle[linewidth=1.5pt](8.2,9.){0.6}
        \end{pspicture*}
    \end{minipage}
\end{bmethode}

\vspace*{-0.6cm}
\section{Développer et factoriser}
\subsection{Développer}

\begin{bdefi}
	\grasInDef{Développer} une expression consiste à transformer un \grasInDef{produit} en une \grasInDef{somme} (ou une différence).
\end{bdefi}

\begin{bprop}
	
	\grasInProp{Distributivité et double distributivité}
	
	Pour tous nombres réels $k$, $a$, $b$, $c$ et $d$ \quad on a :
	
	$$k\times (a+b)=ka+kb$$
	$$(a+b)(c+d)=$$
	
\end{bprop}
\begin{bmethode}
	Quand une parenthèse est précédée d'un signe moins, on développe en multipliant par -1, c'est à dire que l'on change tous les signes à l'intérieur de la parenthèse. Sinon on ne change rien. 
\end{bmethode}


\begin{bexemple}{}
	
	Simplifier :
	$A=\left(-6 x-4\right)-\left(5 x-4\right)=-6x-4-5x+4=-11x$ \\
	Développer : 
	
	$B=(4x- 5)(7-3x)=4x\times 7+4x\times (-3x)-5\times 7-5\times (-3x)=28x-12x^2-35+15x$
	
	$B=-12x^2+43x-35$
\end{bexemple}

\subsection{Factoriser}

\begin{bdefi}
	Factoriser, c'est transformer une \grasInDef{somme (ou une différence)} en un \grasInDef{produit}.\\
\end{bdefi}

\begin{bmethodet}{Factoriser}
	\begin{multicols}{2}
		$A=(2x+3)(4x+1)-(2x+3)(x+2)$\vspace{1cm}\\
		$A=(2x+3)\lbrack......$\vspace{1cm}\\
		$A=(2x+3)\lbrack(4x+1)-(x+2)\rbrack$\\
		\vfill\columnbreak
		\begin{enumMet}
			\item Repérer le facteur commun\\
			Le facteur commun est $(2x+3)$. On peut le souligner dans l'expression de départ.
			\item L'écrire devant et \underline{ouvrir un crochet }
			\item Dans le crochet,on met tout ce qui n'est pas le facteur commun (tout ce qui n'est pas souligné).
		\end{enumMet}
	\end{multicols}
\end{bmethodet}

\begin{bexemple}{}
	Factoriser $B=(x-1)(3x+4)-(x-1)(x+3)$\\
	$B=(x-1)\left[(3x+4)-(x+3)\right]=(x-1)(3x+4-x-3)=(x-1)(2x+1)$
\end{bexemple}

\subsection{Identités remarquable}

\begin{bprop}
	
	\grasInProp{Identités remarquables :}
	
	\noindent Pour tous nombres réels $a$ et $b$  on a :
	%\tcbbreak
	$$\phantom{X} \left(a + b \right)^2 \phantom{Xx} =  a^2 + 2ab +b^2$$
	$$\phantom{X} \left(a - b \right)^2 \phantom{Xx} =  a^2 - 2ab +b^2$$
	$$\left(a + b \right)\left( a- b\right) = a^2 - b^2$$
\end{bprop}

\begin{bexemple}{}
	Factoriser les expressions $P$ et $Q$ ci-dessous.
	\vspace{-0.3cm}
	\setlength{\columnseprule}{0.4pt} 
	\begin{multicols}{2}  
		$P=(x-3)^2-49$ \\
		$P=(x-3)^2-7^2$
		$P=(x-3-7)(x-3+7)$\\
		$P=(x-10)(x+4)$
		
		\columnbreak
		$Q=4x^2-12x+9$ \\
		$Q=(2x)^2-2\times 2x\times 3+3^2$\\○
		$Q=(2x-3)^2$
	\end{multicols} 
\end{bexemple}

%\eject

\section{Calculs avec des fractions}
\subsection{Ajouter ou soustraire des fractions}
\begin{bmethode}
	Pour ajouter ou soustraire des fractions, on les met sous le même dénominateur. \\
	On multiplie le numérateur de la 1ère fraction par le dénominateur de la 2nde.\\
	De la même manière, on multiplie le numérateur de la 2nde par le numérateur de la 1ère.\\
	Le tout se met sous le même dénominateur : le produit des 2 dénominateurs.
\end{bmethode}
\begin{bexemple}
	$\dfrac{1-3x}{x+4}-\dfrac{1+x}{x+2}=\dfrac{(1-3x)\times (x+2)-(1+x) \times (x+4)}{(x+4)(x+2)}$\\
	\phantom{$\dfrac{1-3x+4}{x+1}-\dfrac{1+x}{x+2}$}$=\dfrac{x+2-3x^{2}-3\times 2-(x+4+x^{2}+4x)}{(x+4)(x+2)}$\vspace{0.2cm}\\
	\phantom{$\dfrac{1-3x+4}{x+1}-\dfrac{1+x}{x+2}$}$=\dfrac{x+2-3x^{2}-6-5x-4-x^{2}}{(x+4)(x+2)}$\vspace{0.2cm}\\
	\phantom{$\dfrac{1-3x+4}{x+1}-\dfrac{1+x}{x+2}$}$=\dfrac{-4x-10-4x^{2}}{(x+4)(x+2)}$\vspace{0.2cm}\\
\end{bexemple}

\subsection{Multiplier ou diviser des fractions}
\begin{bmethode}
	Pour multiplier des fractions, on multiplie les numérateurs entre eux, et les dénominateurs entre eux.\\
	Pour diviser des fractions, on multiplie par l'inverse.
\end{bmethode}

\begin{bexemple}
	Simplifier les expressions $A$ et $B$ ci-dessous.
	\vspace{-0.3cm}
	\setlength{\columnseprule}{0.4pt} 
	\begin{multicols}{2}  
		$A=\dfrac{x-3}{3}\times \dfrac{6}{5}$ \\
		$A=\dfrac{(x-3)\times 6}{3\times 5}$\\
		$A=\dfrac{2(x-3)}{5}$
		
		
		\columnbreak
		$B=\dfrac{2}{x-5} \div \dfrac{2}{x+3}$ \\
		$B=\dfrac{2}{x-5} \times \dfrac{x+3}{2}$\\
		$B=\dfrac{2\times (x+3)}{(x-5)\times 2}$\\
		$B=\dfrac{x+3}{x-5}$
	\end{multicols} 
\end{bexemple}

\section{Résolution d'équations}

\subsection{Equation à produit nul}
\begin{bprop}
	Si un produit de facteurs est nul, alors l'un au moins de ses facteurs est nul.
\end{bprop}
\begin{bmethodet}{Se ramener à une équation produit nul}
	\renewcommand{\labelitemi}{\tiny$\bullet$}
	\begin{itemize}
		\item On écrit tous les termes à gauche de l'équation
		\item On FACTORISE $\rightarrow$ soit grâce à un facteur commun\\
		\phantom { On FACTORISE}$\rightarrow$ soit grâce à un identité remarquable
		\item On applique le produit nul.
		
	\end{itemize}
\end{bmethodet}

\begin{bexemple}{}
	Résoudre dans $\R$ l'équation $(2x-1)^2-25= 0$
	
	$(2x-1)^2-25= 0 \iff (2x-1)^2-5^2=0$
	
	$\phantom{(2x-1)^2-25= 0 } \iff (2x-1-5)(2x-1+5)=0$
	
	$\phantom{(2x-1)^2-25= 0 } \iff (2x-6)(2x+4)=0$
	
	EPN : Si un produit de facteurs est nul, alors l'un au moins de ses facteurs est nul.
	
	$\phantom{(2x-1)^2-25= 0 } \iff (2x-6)=0$ ou $(2x+4)=0$
	
	$\phantom{(2x-1)^2-25= 0 } \iff 2x=6$ ou $2x=-4$
	
	$\phantom{(2x-1)^2-25= 0 } \iff x=3$ ou $x=-2$
	$$S=\left\{-2\ ;\ 3\right\}$$
	
\end{bexemple}

\subsection{Equations du type $x^2=a$}

\begin{bmethode}
	Les solutions de l'équation sont $x=\sqrt{a}$ et $x=-\sqrt{a}$ si a est positif.
	
	Il n'y a pas de solutions si a est négatif (un carré est toujours positif).
\end{bmethode}

\begin{bexemple} 
	$x^{2}=7$
	$$S=\left\{-\sqrt{7}\ ;\ \sqrt{7}\right\}$$  
\end{bexemple}

\section {Puissances}
\begin{bdefi} 
	Pour tout nombre réel $a$ et pour tout entier naturel $n$, on a : \\
	$a^{0}=1$  \quad \quad \quad \quad \quad $a^1=a$ \quad \quad \quad \quad \quad \quad $a^{n}=a \times a \times a \ldots \times a$ ($n$ termes) \quad \quad \quad \quad \quad $a^{-n}=\dfrac{1}{a^{n}}$
\end{bdefi}



\begin{bexemple}{}
	\begin{multicols}{5}
		$3^0=1$\par
		$(-5)^1=-5$ \par 
		$(-7)^0=1$ \par
		$7^{-2}=\dfrac{1}{7^2}=\dfrac{1}{49}$ \par
		$2^3=2\times 2\times 2=8$ \par 
	\end{multicols}
	
\end{bexemple}

\begin{proprietes}{Si les nombres élevés à la puissance sont les mêmes : (ici a)}
	\begin{multicols}{3}
		$a^{n+m}=a^{n} \times a^{m}$ \par
		$a^{n-m}=\dfrac{a^{n}}{a^{m}}$ \par
		${(a^{n})}^m=a^{n \times m}$
	\end{multicols}
\end{proprietes}

\begin{bexemple}{}
	\begin{multicols}{2}
		$7^{5} \times 7^{9}=7^{5+9}=7^{14}$ \par
		$5^{3} \times 5^{2} \times 5^{-5}=5^{3+2-5}=3^0=1$ \par
		$\dfrac{6^{8}}{6^{7}}=6^{8-7}=6^1=6$ \par
		$\left(3^{-4}\right)^{5}=3^{-4\times 5}=3^{-20}$ \par
	\end{multicols}
\end{bexemple}


\begin{proprietes}{Si les puissances sont les mêmes : (ici n)}
	\begin{multicols}{2}
		$a^{n} \times b^{n}=(a \times b)^{n}$ \par
		$\dfrac{a^{n}}{b^{n}}=\left(\dfrac{a}{b}\right)^{n}$ \par
		
	\end{multicols}
\end{proprietes}

\begin{bexemple}{}
	\begin{multicols}{2}
		$7^{5} \times 4^{5}=(7\times 5)^5=35^5$ \par
		$\dfrac{6^{8}}{4^{8}}=\left(\dfrac{6}{8}\right)^8=\left(\dfrac{3}{4}\right)^8$ \par
	\end{multicols}
\end{bexemple}

\begin{brmqs}
	
	\begin{itemRem}
		\item Si les nombres élevés à la puissance sont différents et si les puissances sont différentes, il n’existe pas de règle de calcul : 
		
		\begin{multicols}{2}
			$7^{5} \times 4^{3}=??$ \par
			$\dfrac{5^{8}}{4^{5}}=??$ \par
		\end{multicols}
		
		\item Il n'y a pas de règle pour additionner ou soustraire des puissances : il faut factoriser : 
		
		$2^{31}-2^{30}= 2^{30+1}-2^{30}=2\times 2^{30}-1\times 2^{30}=2^{30}(2-1)=2^{30}$
		
	\end{itemRem}
	
\end{brmqs}

\subsection*{Puissances et écriture scientifique}
\begin{bdefi}
	En notation scientifique, un nombre s'écrit sous la forme $a \times 10^n$ où $a$ est un nombre décimal supérieur ou égal à 1 et strictement inférieur à 10, et $n$ est un entier relatif.
\end{bdefi}


\begin{bexemple}
	Ecrire en notation scientifique : \\
	\begin{multicols}{2}
		\begin{enumEx}
		\item $546,7 \times 10^{9}=5,467\times 10^{11}$
		\item $0,045 \times 10^{-3}=4,5\times 10^{-5}$
	\end{enumEx}
	\end{multicols}
\end{bexemple}



\end{document}